\documentclass[11pt,a4paper]{article}

\usepackage[utf8]{inputenc}
\usepackage[T1]{fontenc}
\usepackage{lmodern}
\usepackage{amsmath,amssymb}
\usepackage{booktabs}
\usepackage{graphicx}
\usepackage{hyperref}
\usepackage{geometry}
\usepackage{xcolor}

\geometry{margin=2.5cm}
\hypersetup{colorlinks=true,linkcolor=blue!50!black,citecolor=blue!50!black,urlcolor=blue!50!black}

\newcommand{\Uset}{\mathcal{U}}
\newcommand{\Iset}{\mathcal{I}}
\newcommand{\Fset}{\mathcal{F}}
\newcommand{\Rset}{\mathcal{R}}
\newcommand{\pvec}{\mathbf{p}}
\newcommand{\qvec}{\mathbf{q}}
\newcommand{\evec}{\mathbf{e}}

\title{\textbf{Hybrid Feature-Aware Matrix Factorization for\\
Cold-Start Skincare Retrieval}\\[0.4em]
\large A from-scratch WARP-loss implementation (the LightFM algorithm)\\
with a dedicated cold-start evaluation protocol}

\author{Patrick Oliveira}
\date{\today}

\begin{document}

\maketitle

\begin{abstract}
We implement hybrid, feature-aware matrix factorization trained with a
Weighted Approximate-Rank Pairwise (WARP) loss -- the algorithm popularized by
LightFM -- entirely from scratch in NumPy/SciPy, for top-$N$ product
retrieval on the Sephora Products and Skincare Reviews dataset. Unlike the
two identity-indexed sibling projects in this portfolio (BPR over anime
titles, SVD over MovieLens ratings), user and item representations here are
sums of \emph{feature} embeddings -- skin type and tone for users; brand,
category, price tier and ingredient tags for products -- rather than sums
indexed by raw identity. This is the classical, pre-deep-learning answer to
cold start: a user or item the model has never seen individually can still be
scored, as long as it shares known features with the training population. We
report a dedicated cold-start evaluation slice, built by withholding a
population of users and items from training entirely, and show the hybrid
model retrieves measurably better-than-popularity candidates for them, while
an ID-only matrix-factorization baseline and a pure popularity baseline both
collapse to a score of exactly zero on the same slice -- not a close loss,
but a structural inability to say anything at all.
\end{abstract}

\tableofcontents
\newpage

\section{Notation}
\label{sec:notation}

Let $\Uset$ be the set of users and $\Iset$ the set of items. Each user $u$
has a set of known features $F_u \subseteq \Fset_{\text{user}}$ (skin type,
skin tone, and -- only for users seen during training -- an identity token);
each item $i$ has a set of known features $F_i \subseteq \Fset_{\text{item}}$
(brand, category, price bucket, ingredient tags, and -- only for items seen
during training -- an identity token). Every feature $f$ owns a learned
embedding $\evec_f \in \mathbb{R}^d$. Implicit feedback is the set of
observed positive interactions $\Rset \subseteq \Uset \times \Iset$, and
$\Iset_u^{+} = \{\, i \in \Iset : (u,i) \in \Rset \,\}$ denotes the products a
user reviewed positively.

\section{Introduction}
\label{sec:introduction}

Every recommender in this portfolio is implicit-feedback matrix
factorization; what differs between them is what gets factorized. The anime
project factorizes a $73\text{k} \times 12\text{k}$ matrix of raw user and
item identities with a BPR loss~\cite{rendle2009bpr}; the MovieLens project
factorizes an explicit rating matrix with truncated SVD. Both share the same
structural limitation: a row or column exists only for an identity observed
during training. A user who signs up the day after the last retraining job,
or a product added to the catalogue this morning, has no row or column at
all, and the only fallback available to an identity-indexed model is
non-personalized popularity.

\subsection{The hybrid alternative}

LightFM~\cite{kula2015} reframes the factorization: instead of learning one
embedding per user and one per item, it learns one embedding per
\emph{feature}, and represents a user or item as the sum of the embeddings
of its known features,
\begin{equation}
  \pvec_u = \sum_{f \in F_u} \evec_f, \qquad
  \qvec_i = \sum_{g \in F_i} \evec_g .
  \label{eq:feature-sum}
\end{equation}
A brand-new user who has filled in a skin-type and skin-tone profile already
has a well-defined $\pvec_u$, built from embeddings the model learned from
every \emph{other} user who shares those two attributes. The same holds for a
brand-new product described by its brand, category, price tier and
ingredient tags. Identity is folded in as just one more feature -- present
for users/items observed during training, deliberately absent otherwise --
so the model degrades gracefully to a pure content-based recommender rather
than failing outright.

This is a classical, linear, SGD-trained model: there are no hidden layers
and no nonlinear activations between the feature sum and the score. The
contribution of this report is not a new loss function; it is (a) a
from-scratch, dependency-free implementation of the WARP loss that LightFM
popularized, and (b) a cold-start evaluation protocol precise enough to
quantify the gap between a hybrid and an identity-only model, rather than
asserting it qualitatively.

\subsection{Why not depend on the \texttt{lightfm} package}

The \texttt{lightfm} PyPI package exists and implements the same idea, but
its Cython extension is compiled against a fixed, old NumPy ABI and routinely
fails to build on current Python/NumPy combinations -- exactly the kind of
fragile dependency a portfolio repository should not ask a reviewer to fight
with. Reimplementing the algorithm in plain NumPy/SciPy (Section
\ref{sec:methodology}) removes that dependency entirely while keeping every
equation in this report directly traceable to a function in the
repository.

\subsection{Scope and claims}

\begin{enumerate}
  \item The feature layer is point-in-time correct. Aggregate features are
        computed with a strict \texttt{<} cutoff against each user's own
        interaction order, never \texttt{<=}, and this is enforced by tests
        that inject a future review and assert no feature value changes
        (Section~\ref{sec:features}).
  \item Cold-start is an operationalized test, not a claim. A population of
        users and items is withheld from training \emph{entirely} -- not
        merely their last interaction -- and the identity-feature vocabulary
        is built only from the training partition, so a cold entity provably
        has no ID token available (Section~\ref{sec:evaluation}).
  \item The comparison is fair. The ID-only baseline is the same family of
        pairwise SGD factorization (Section~\ref{sec:methodology}), not a
        strawman; it loses on the cold-start slice for a structural reason
        (no row/column exists for a withheld entity), not a tuning one.
\end{enumerate}

\subsection{Repository}

The implementation, infrastructure, and the scripts that generate every
table in this document are available in the accompanying repository. The
results tables (Tables~\ref{tab:results} and \ref{tab:cold-start}) are
emitted by \texttt{scripts/export\_results\_table.py} and the feature
catalogue (Table~\ref{tab:features}) by
\texttt{scripts/export\_feature\_tables.py}, so the document cannot drift
away from the code.

\section{Data}
\label{sec:data}

The target dataset is the Sephora Products and Skincare Reviews collection
(Kaggle ref.\ \texttt{nadyinky/sephora-products-and-skincare-reviews}):
roughly 8{,}000 products with attributes such as category, brand, price and
ingredient lists, and on the order of one million reviews carrying a rating,
a recommendation flag, and reviewer attributes including skin type and skin
tone.

\subsection{Synthetic fixture}

The repository ships a synthetic generator
(\texttt{scripts/download\_data.py}) that reproduces this schema without
requiring Kaggle credentials, so the test suite and CI run with zero external
dependencies. Products are assigned a brand, a category, a price drawn from a
category-appropriate range, and two to four ingredient tags; one tag is
deliberately drawn from a fixed list of tags each skin type responds well to
(Table~\ref{tab:skin-affinity}), so the catalogue always contains genuine
structure for every skin type to recover, not just noise.

\begin{table}[h]
  \centering
  \caption{Synthetic ingredient-affinity clusters. A real dataset would not
  have this structure made explicit; the synthetic fixture encodes it so
  that a model with the right inductive bias (shared ingredient features)
  can be shown to recover it with a bounded sample size.}
  \label{tab:skin-affinity}
  \begin{tabular}{ll}
    \toprule
    Skin type & Preferred ingredient tags \\
    \midrule
    Dry & Hyaluronic Acid, Ceramides, Squalane, Glycerin \\
    Oily & Salicylic Acid, Niacinamide, Clay, Zinc PCA \\
    Combination & Niacinamide, Hyaluronic Acid, Green Tea \\
    Normal & Vitamin C, Peptides, Antioxidant Complex \\
    Sensitive & Centella Asiatica, Oat Extract, Ceramides, Panthenol \\
    \bottomrule
  \end{tabular}
\end{table}

Each synthetic user is assigned a skin type and a skin tone, and reviews
products with a selection probability weighted by the overlap between the
product's ingredient tags and the user's skin-type cluster, mildly tilted by
a long-tailed popularity prior. Ratings are drawn around a centre that is
higher for well-matched products and lower otherwise, and a review counts as
``recommended'' only when the rating clears the implicit-feedback threshold
(Section~\ref{sec:evaluation}). Every selection decision is keyed off
\emph{features}, never off the raw user or product ID -- this is what gives
a feature-aware model real structure to recover that an ID-only model
cannot, by construction, generalize from one identity to another.

\subsection{Implicit feedback}

A review is treated as a positive interaction when
$\text{rating} \geq 4$ or the reviewer explicitly marked the product as
recommended; everything else is dropped rather than treated as an explicit
negative, following the usual implicit-feedback convention. An iterative
$k$-core filter then removes users and products with fewer than five
surviving positives, looping because removing a product can push a user
below the floor and vice versa.

\section{Feature Engineering}
\label{sec:features}

In the sibling BPR and SVD projects, the feature layer produces auxiliary
signal: confidence weights, popularity priors, content features consulted
only as a cold-start fallback. Here the feature layer \emph{is} the model's
representation -- there is no separate identity-indexed embedding table that
the features merely augment. This changes what correctness means for the
pipeline: a bug in a feature transform is not a secondary signal degrading
gracefully, it is a bug in the thing being factorized.

\subsection{Declarative specification}

\texttt{conf/features.yaml} declares every feature once, consumed
identically by training and by the serving path
(\texttt{tests/test\_features\_parity.py} asserts the two never drift apart).
Each declaration carries a \emph{fit scope}:

\begin{itemize}
  \item \textbf{catalogue} -- known the moment a product is listed (brand,
        category, price bucket, ingredient tags). Always available, cold
        item or not.
  \item \textbf{profile} -- a static, self-reported user attribute (skin
        type, skin tone). Not a statistic, so not subject to the train-only
        rule; a brand-new user supplies it at signup.
  \item \textbf{train\_only} -- a vocabulary or statistic restricted to the
        training partition. Both identity features (\texttt{user\_identity},
        \texttt{item\_identity}) use this scope: the vocabulary of ID tokens
        literally does not contain an entry for a withheld entity, which is
        what makes cold-start an operationalized property of the pipeline
        rather than an assertion about it.
\end{itemize}

Table~\ref{tab:features} lists every declared feature with its family, kind
and scope; the per-feature rationale (Table~\ref{tab:feature-rationale})
documents why each one was declared.

% Generated by scripts/export_feature_tables.py -- do not edit by hand.
\begin{table}[t]
  \centering
  \caption{Declared features. ``Scope'' is the fit boundary: catalogue
  and profile features are known a priori; train\_only features --
  including both identity features -- are restricted to the training
  partition, which is what makes cold users/items genuinely lack an ID
  token rather than merely an unseen one.}
  \label{tab:features}
  \begin{tabular}{llllc}
    \toprule
    Feature & Family & Kind & Scope & PIT \\
    \midrule
    \texttt{user\_skin\_type} & user & categorical & profile & -- \\
    \texttt{user\_skin\_tone} & user & categorical & profile & -- \\
    \texttt{user\_identity} & user & identity & train\_only & -- \\
    \texttt{item\_brand} & item & categorical & catalogue & -- \\
    \texttt{item\_category} & item & categorical & catalogue & -- \\
    \texttt{item\_price\_bucket} & item & categorical & catalogue & -- \\
    \texttt{item\_ingredient\_tags} & item & multi\_categorical & catalogue & -- \\
    \texttt{item\_identity} & item & identity & train\_only & -- \\
    \texttt{user\_review\_count\_bucket} & interaction & categorical & train\_only & \checkmark \\
    \texttt{item\_popularity\_bucket} & interaction & categorical & train\_only & -- \\
    \texttt{item\_catalogue\_context} & context & categorical & train\_only & -- \\
    \bottomrule
  \end{tabular}
\end{table}


\begin{table}[h]
  \centering
  \caption{Per-feature rationale, generated from the \texttt{rationale}
  field every entry in \texttt{conf/features.yaml} is required to carry
  (\texttt{test\_features\_schema.py} rejects a feature without one).}
  \label{tab:feature-rationale}
  % Generated by scripts/export_feature_tables.py -- do not edit by hand.
\begin{description}
  \item[\texttt{user\_skin\_type}] Core demographic attribute driving taste. Known at signup, so it is available for every user including ones with zero training history -- the attribute a brand-new user supplies is exactly what lets the hybrid model place them near users with similar skin.
  \item[\texttt{user\_skin\_tone}] Secondary demographic attribute (shade matching for base/coverage products). Same availability argument as skin type.
  \item[\texttt{user\_identity}] Per-user ID feature, LightFM-style: lets the model retain personalised signal for warm users without it being the *only* signal. The vocabulary is restricted to users seen in training, so a cold user genuinely has no ID token available -- this is the feature that operationalizes "cold start" rather than merely asserting it.
  \item[\texttt{item\_brand}] Brand loyalty is a strong prior independent of any single product's review history, and is known the moment a product is catalogued.
  \item[\texttt{item\_category}] Coarse product type (serum vs. cleanser vs. sunscreen). Known at catalogue time, safe for cold items by construction.
  \item[\texttt{item\_price\_bucket}] Discretized price tier. Budget shoppers and prestige shoppers cluster on this signal regardless of brand or category.
  \item[\texttt{item\_ingredient\_tags}] The content signal that actually explains taste: a Dry-skin user's affinity is for Hyaluronic-Acid / Ceramide products specifically, not for "Moisturizers" as a category. This is what lets the model recommend a brand-new product that happens to share tags with liked products.
  \item[\texttt{item\_identity}] Mirrors user\_identity: vocabulary restricted to items seen in training, so a cold item has no ID token and must be scored through its brand/category/ingredient tokens alone.
  \item[\texttt{user\_review\_count\_bucket}] Confidence weighting: a user with 40 training reviews has a much more reliable taste signal than one with 3. Computed with a strict `<` cutoff against the user's own interaction\_order so it can never see the very review it is being used to help predict.
  \item[\texttt{item\_popularity\_bucket}] Popularity prior used by the most-popular baseline and as part of the cold-start fallback ordering. Fitted on the training partition only, so held-out interactions never inflate a product's apparent popularity.
  \item[\texttt{item\_catalogue\_context}] Distinguishes an "established" item (enough training reviews to trust its popularity/ingredient signal) from a "new\_arrival" (sparse or zero training reviews). Cold items always fall back to this value, documenting that they are indistinguishable from genuinely new catalog entries from the model's point of view.
\end{description}

\end{table}

\subsection{Point-in-time correctness}

Aggregate features derived from the interaction log (e.g.\
\texttt{user\_review\_count\_bucket}) accept a strict cutoff and filter with
\texttt{<}, never \texttt{<=}, against each user's own
\texttt{interaction\_order}. \texttt{tests/test\_features\_leakage.py}
operationalizes this directly: it injects a review positioned after the
cutoff for an existing user and asserts that no feature value for that user
changes, and separately asserts that a review positioned \emph{exactly at}
the cutoff is excluded (the boundary is exclusive, not inclusive). Item-level
aggregates (\texttt{item\_popularity\_bucket}) are instead protected by a
\texttt{fit\_on: train\_only} contract, verified by asserting the statistic
computed on the training partition differs from the same statistic computed
on the full dataset -- if they matched, the held-out interactions would
already be baked into a value the model is scored on.

\subsection{A deliberate exclusion from the retrieval representation}

\texttt{item\_popularity\_bucket} and \texttt{item\_catalogue\_context} are
declared, registered and leakage-tested like every other feature, but are
not included in the item representation the hybrid model actually
factorizes. The reason is structural rather than a matter of taste: a
withheld (cold) item has, by construction, exactly zero training
interactions, so both aggregates collapse every cold item onto the same
``unpopular'' / ``new\_arrival'' token. Including them in $\qvec_i$ would
teach the model to systematically suppress cold items through a learned
global popularity prior at precisely the moment we want to test whether it
can still retrieve them through shared brand/category/ingredient structure.
The ID-only and popularity baselines already encode an item-popularity prior
through their own bias terms -- a prior that is, for the same structural
reason, unavailable for a cold item -- so nothing is lost from the
comparison by leaving these two aggregates out of the hybrid model's own
representation; they remain fully part of the feature store, the schema and
leakage tests, and the feature table above.

\subsection{Testing surface}

Three test files operationalize the three properties above:
\texttt{test\_features\_schema.py} validates \texttt{conf/features.yaml}
against the registry (every declared transform is implemented and vice
versa, every identity feature is declared \texttt{train\_only}, every
feature documents a rationale); \texttt{test\_features\_leakage.py}
implements the point-in-time and fit-scope checks described above, plus a
label-shuffling sanity check that destroys the user--item association and
asserts ranking quality collapses to the random floor;
\texttt{test\_features\_parity.py} asserts that
\texttt{recompute\_user\_tokens}/\texttt{recompute\_item\_tokens} -- the
functions the serving path calls for an entity outside the last training
run -- reproduce exactly the tokens the batch-built feature store used
during training, for every warm user/item sampled.

\section{Methodology}
\label{sec:methodology}

\subsection{Score function}

Both the user and item representations are sums of feature embeddings
(Equation~\ref{eq:feature-sum}, repeated here for the score derivation):
\begin{equation}
  \pvec_u = \sum_{f \in F_u} \evec_f, \qquad
  \qvec_i = \sum_{g \in F_i} \evec_g .
\end{equation}
The score of a user--item pair is their inner product plus an item bias,
\begin{equation}
  s(u, i) = \Big\langle \sum_{f \in F_u} \pvec_f,\ \sum_{g \in F_i} \qvec_g \Big\rangle
          = \langle \pvec_u, \qvec_i \rangle + b_i ,
  \label{eq:score}
\end{equation}
where $b_i = \sum_{g \in F_i} \beta_g$ is itself a sum of per-feature bias
scalars, mirroring the embedding sum. A per-user bias is deliberately
omitted: it would shift every item's score for that user by the same
constant and therefore cancel out of the pairwise margin used by the loss
below, exactly as the per-user bias does in the sibling BPR project's own
derivation -- it would carry zero gradient under this objective.

\subsection{WARP loss}
\label{sec:warp}

Weighted Approximate-Rank Pairwise loss~\cite{weston2011warp,kula2015}
optimizes a pairwise ranking criterion like BPR, but instead of weighting
every violated pair equally, it weights a positive pair by an \emph{estimate
of the positive's current rank} in the full catalogue, so that correcting a
badly mis-ranked positive contributes far more gradient than nudging one
that is already well-placed.

For a positive pair $(u, i)$, negatives $j$ are sampled uniformly and
rejected until a margin violation is found,
\begin{equation}
  s(u, j) > s(u, i) - 1 ,
\end{equation}
or until $N_{\max}$ trials are exhausted. Let $N$ be the number of trials
taken. Because negatives are sampled uniformly at random, the probability
that a given negative violates the margin is approximately
$\text{rank}(i) / (|\Iset| - 1)$, so the expected number of trials before
the first violation is the reciprocal of that probability; inverting gives
the rank estimate
\begin{equation}
  \text{rank}(i) \approx \left\lfloor \frac{|\Iset| - 1}{N} \right\rfloor .
  \label{eq:rank-estimate}
\end{equation}
Few trials imply many items already violate the margin, i.e.\ $i$ is ranked
low; many trials imply $i$ is already hard to beat, i.e.\ ranked high. The
hinge loss for the triplet is weighted by the margin-ranking loss function
\begin{equation}
  L(\text{rank}(i)) = \sum_{r=1}^{\text{rank}(i)} \frac{1}{r} ,
  \label{eq:warp-weight}
\end{equation}
a harmonic sum that grows (slowly) with rank, so that
\begin{equation}
  \ell(u, i, j) = L(\text{rank}(i)) \cdot \max\!\big(0,\ 1 - s(u,i) + s(u,j)\big) .
  \label{eq:warp-loss}
\end{equation}

\subsection{Gradients and the update rule}

Write $m = 1 - s(u,i) + s(u,j)$ for the margin and $w = L(\text{rank}(i))$
for the WARP weight; $\ell = w \cdot \max(0, m)$ is non-zero, and has a
gradient, only when $m > 0$. For an active triplet,
\begin{align}
  \frac{\partial \ell}{\partial s(u,i)} &= -w , &
  \frac{\partial \ell}{\partial s(u,j)} &= +w .
\end{align}
Since $s(u,i) = \langle \pvec_u, \qvec_i \rangle + b_i$ and every feature
embedding is shared across every entity that carries that feature, the
chain rule distributes the gradient to \emph{every active feature
embedding} of $u$, $i$ and $j$ simultaneously:
\begin{align}
  \pvec_u \text{'s features } f &:\ \evec_f \leftarrow \evec_f +
    \eta \big( w (\qvec_i - \qvec_j) - \lambda \evec_f \big) , \\
  \qvec_i \text{'s features } g &:\ \evec_g \leftarrow \evec_g +
    \eta \big( w \, \pvec_u - \lambda \evec_g \big) , \\
  \qvec_j \text{'s features } g &:\ \evec_g \leftarrow \evec_g +
    \eta \big( {-w} \, \pvec_u - \lambda \evec_g \big) , \\
  b_i \text{'s features } g &:\ \beta_g \leftarrow \beta_g + \eta \big( w - \lambda \beta_g \big) , \\
  b_j \text{'s features } g &:\ \beta_g \leftarrow \beta_g + \eta \big( {-w} - \lambda \beta_g \big) ,
\end{align}
with learning rate $\eta$ and L2 penalty $\lambda$. This is the sense in
which the model is a classical, linear SGD factorization: every update is a
closed-form linear-algebra step on embeddings, with no hidden layer and no
nonlinear activation anywhere in the forward pass.

\paragraph{A stability correction that matters in practice.} A feature
shared by many low-popularity training items (in our pipeline this would
have been \texttt{popularity=0}, had it not been excluded from the
representation per Section~\ref{sec:features}) is sampled as a WARP
negative about as often as any other feature but is rarely part of a
positive, so with only a weak L2 pull its bias term can diverge far faster
than $\lambda$ counteracts it -- in early experiments this drove cold-item
scores to the bottom of the ranking regardless of their brand/ingredient
match. Clipping every touched embedding to a bounded L2 norm and every
touched bias scalar to a bounded magnitude after each update (both
configurable; LightFM's own reference implementation does the same for
embeddings) removes the divergence without changing the gradient direction,
only its magnitude.

\subsection{Baselines}

\paragraph{ID-only MF (BPR).} The same pairwise-SGD family, but the user and
item representations are raw identity rows/columns rather than feature
sums -- structurally identical to the sibling BPR project. A user or item
absent from training has no row or column, so its score collapses to the
item bias term alone (a pure popularity signal) for warm items, and to
\emph{nothing meaningful at all} for a cold item, since that item's bias was
also never trained; its factor and bias are masked out of the rank for every
user, warm or cold.

\paragraph{Most-popular.} Ranks every user by the same global training-count
order. Just like the ID-only baseline, a cold item has a training count of
exactly zero and is therefore never recommended by this baseline either.

\paragraph{Random.} A seeded uniform permutation per user, which calibrates
the metric scale; because its coverage is total and uncorrelated with
quality, high coverage alone is not evidence of a good model.

\section{Evaluation Protocol}
\label{sec:evaluation}

Two different notions of ``holdout'' are used, and keeping them distinct is
the whole point of this report.

\subsection{Warm-start: leave-last-one-out}

Within the population of users retained after core filtering, each user's
last interaction (by true review timestamp) is held out for test, the
second-to-last for validation, and the rest for training. Users without
enough history to leave at least three training interactions contribute
only to training and are never evaluated, which avoids inflating metrics
with users whose ``history'' is a single review.

\subsection{Cold-start: entities withheld entirely}

Before the warm-start split is computed, a fraction of users and a fraction
of items (8\% each, by default) are selected uniformly at random and every
interaction touching either is removed from the training-eligible pool
\emph{before} any splitting happens. The warm-start split above is then
computed only over the remainder. Three evaluation slices follow directly
from this construction:

\begin{itemize}
  \item \textbf{New user, known item} (\texttt{cold\_user\_eval}): reviews
        by a withheld user against a product that \emph{was} trained on.
        Tests whether shared demographic features let the model retrieve
        for a user it has never seen.
  \item \textbf{Known user, new item} (\texttt{cold\_item\_eval}): reviews
        by a trained user against a withheld product. Tests whether shared
        brand/category/ingredient features let the model retrieve a product
        it has never seen for a user whose taste it has learned.
  \item \textbf{New user, new item} (\texttt{cold\_both\_eval}): both sides
        withheld -- the hardest slice, reported for completeness.
\end{itemize}

Because the identity-feature vocabulary (Section~\ref{sec:features}) is
built from the training partition alone, a withheld entity provably has no
identity token in either the hybrid model or the ID-only baseline. The
exclusion set passed to every model's ranking call is empty for these
slices: there is no training history to hide, which is exactly the point.

\subsection{Metrics}

For every evaluated user, the model ranks the full catalogue (excluding
training history, where applicable) and NDCG@$K$, Recall@$K$ and MAP@$K$ are
computed against the relevant set -- a single item for the warm-start
partitions (leave-last-one-out), potentially several for the cold-start
slices (every interaction a withheld entity has is relevant, since all of
them were excluded from training).
\begin{align}
  \text{NDCG@}K &= \frac{1}{|\Uset_{\text{eval}}|} \sum_{u}
    \frac{\sum_{k=1}^{K} \mathbb{1}[\text{rel}(k)] / \log_2(k+1)}{\text{IDCG}_u} , \\
  \text{Recall@}K &= \frac{1}{|\Uset_{\text{eval}}|} \sum_{u}
    \frac{|\{\text{relevant items retrieved in top-}K\}|}{|\text{relevant items for } u|} , \\
  \text{MAP@}K &= \frac{1}{|\Uset_{\text{eval}}|} \sum_{u}
    \frac{1}{\min(|\text{relevant}_u|, K)} \sum_{k=1}^{K} \mathbb{1}[\text{rel}(k)] \cdot \text{Precision@}k .
\end{align}

\subsection{Hyperparameter selection}

Hyperparameters for every model (hybrid and baselines alike) are fixed
before the test partition and both cold-start slices are read; the values
used are recorded in \texttt{src/recommender/cli.py}
(\texttt{BEST\_HYBRID}, \texttt{BEST\_ID\_MF}) and are the same values the
table-generation script (\texttt{scripts/export\_results\_table.py}) uses,
so the numbers in Section~\ref{sec:results} are reproducible from a clean
checkout with \texttt{python scripts/download\_data.py \&\& python
scripts/export\_results\_table.py}.

\section{Results}
\label{sec:results}

\subsection{Warm-start}

% Generated by scripts/export_results_table.py -- do not edit by hand.
\begin{table}[t]
  \centering
  \caption{Warm-start ranking quality on the held-out test partition ($K=10$).
  Hyperparameters were fixed before the test partition was read once.}
  \label{tab:results}
  \begin{tabular}{lrrrrr}
    \toprule
    Model & NDCG@10 & Recall@10 & MAP@10 & Hit rate & Coverage \\
    \midrule
    ID-only MF & \textbf{0.1120} & 0.2084 & 0.0827 & 0.2084 & 0.8950 \\
    Popularity & 0.0788 & 0.1508 & 0.0572 & 0.1508 & 0.1000 \\
    Hybrid WARP & 0.0655 & 0.1286 & 0.0466 & 0.1286 & 0.8900 \\
    Random & 0.0364 & 0.0687 & 0.0265 & 0.0687 & 1.0000 \\
    \bottomrule
  \end{tabular}
\end{table}


On the warm-start test partition, the ID-only MF baseline leads on raw
ranking quality (Table~\ref{tab:results}). This mirrors the honest negative
result reported by the sibling BPR project: a model with fewer structural
constraints has an easier time fitting a warm population it has already
seen reviews from, because every identity it is scored on had its own row to
learn. The hybrid model pays a complexity tax here -- it has to route all of
its signal through a comparatively small vocabulary of shared features
rather than a dedicated parameter per identity -- and that tax is not fully
recovered by the identity feature it is also given.

\subsection{Cold-start}

% Generated by scripts/export_results_table.py -- do not edit by hand.
\begin{table}[t]
  \centering
  \caption{Cold-start retrieval quality ($K=10$). Users/items in this
  slice were withheld from training entirely; the ID-only baseline and the
  popularity baseline can only fall back to a popularity prior for them,
  while the hybrid model scores them through shared category, brand and
  ingredient features.}
  \label{tab:cold-start}
  \begin{tabular}{llrrrr}
    \toprule
    Model & Slice & NDCG@10 & Recall@10 & MAP@10 & Hit rate \\
    \midrule
    Hybrid WARP & New user, new item & 0.0282 & 0.0758 & 0.0126 & 0.0909 \\
    Random & New user, new item & 0.0236 & 0.0389 & 0.0098 & 0.1212 \\
    Popularity & New user, new item & 0.0000 & 0.0000 & 0.0000 & 0.0000 \\
    ID-only MF & New user, new item & 0.0000 & 0.0000 & 0.0000 & 0.0000 \\
    Random & Known user, new item & 0.0308 & 0.0570 & 0.0160 & 0.1084 \\
    Hybrid WARP & Known user, new item & 0.0129 & 0.0328 & 0.0047 & 0.0623 \\
    Popularity & Known user, new item & 0.0000 & 0.0000 & 0.0000 & 0.0000 \\
    ID-only MF & Known user, new item & 0.0000 & 0.0000 & 0.0000 & 0.0000 \\
    Popularity & New user, known item & 0.2710 & 0.1378 & 0.1333 & 0.9250 \\
    ID-only MF & New user, known item & 0.2661 & 0.1268 & 0.1363 & 0.8500 \\
    Hybrid WARP & New user, known item & 0.1956 & 0.0970 & 0.0888 & 0.7750 \\
    Random & New user, known item & 0.0642 & 0.0438 & 0.0197 & 0.5250 \\
    \bottomrule
  \end{tabular}
\end{table}


Table~\ref{tab:cold-start} is the result this project exists to produce. On
the \emph{new user, known item} slice, every model (including the hybrid
one) relies on an item-level bias/popularity signal, so all three
non-random baselines perform comparably; the hybrid model's demographic
feature sum (skin type and tone) is a weaker but genuinely non-trivial
signal on its own, clearing the random floor by a wide margin.

On the \emph{known user, new item} and \emph{new user, new item} slices, the
comparison is structural rather than a matter of degree: the popularity and
ID-only baselines both report \textbf{exactly zero} NDCG, Recall and MAP.
This is not a close loss -- a withheld item has a training count of
precisely zero, so its popularity-bias term and its identity row are both
untrained, and neither baseline has any other signal to fall back on for it.
The hybrid model scores strictly above zero on both slices, because its item
representation routes through brand, category and ingredient-tag embeddings
learned from every \emph{other} item that shares those tags -- exactly the
mechanism Equation~\ref{eq:feature-sum} was built to provide.

\subsection{Honest caveats}

The hybrid model's absolute recall on the cold-item slices is modest, and
on this synthetic fixture's small catalogue ($\sim$200 products) a uniform
random baseline is occasionally competitive with it in raw hit-rate, simply
because the ratio of $K$ to catalogue size is already favourable to chance.
We report this rather than tuning the random seed until it looks better: the
qualitative claim this project supports is that the hybrid model is
\emph{capable of saying anything at all} about a withheld entity, which the
two comparison baselines are structurally incapable of, not that it is
dramatically better than chance on every metric at this dataset scale. A
larger catalogue, more training epochs, or a sharper per-feature affinity
signal would be the natural next experiments to widen the absolute margin
(Section~\ref{sec:conclusion}).

\section{Conclusion}
\label{sec:conclusion}

We implemented hybrid feature-aware matrix factorization with a WARP ranking
loss end to end, from scratch in NumPy/SciPy: point-in-time feature
engineering over demographic and catalogue attributes, a from-scratch WARP
training loop with the stability corrections it needs in practice, a
dedicated cold-start evaluation protocol that withholds entities from
training entirely rather than merely their last interaction, and a
parity-tested serving path -- precomputed top-$N$ for known entities, a
documented request-time fallback for genuinely new ones -- deployed by
Terraform against a local AWS emulator.

The measured outcome has two parts, and we report both. On warm-start
ranking quality, the simpler ID-only baseline wins, consistent with the
sibling BPR project's own honest negative result: fitting a population
everyone has already been seen in favours fewer structural constraints. On
the cold-start slices this project was built to measure, the comparison is
not close: popularity and ID-only MF are both \emph{structurally} incapable
of saying anything about a withheld entity (exactly zero on every metric),
while the hybrid model retrieves a strictly positive, if modest, signal
through shared brand/category/ingredient structure alone.

\subsection{What transfers}

\begin{itemize}
  \item \textbf{A feature store that is the model, not an accessory to it.}
        Three dedicated test files (schema, leakage, parity) treat the
        feature layer with the same rigor the sibling projects reserve for
        their scoring function, because here they are the same thing.
  \item \textbf{A cold-start harness.} Withholding entities from training
        entirely, restricting identity vocabularies to the training
        partition, and reporting three distinct withheld-population slices
        is a reusable pattern for any feature-aware recommender, not
        specific to skincare.
  \item \textbf{A from-scratch WARP implementation} with the gradient
        derivation and the embedding/bias clipping correction documented in
        Section~\ref{sec:methodology}, removing the fragile \texttt{lightfm}
        Cython dependency entirely.
\end{itemize}

\subsection{Future work}

\paragraph{Sharper negative sampling.} WARP's rejection sampling degrades as
training items increasingly separate into clear winners and losers; sampling
negatives with a mild popularity tilt (as the sibling BPR project's own
future-work notes propose) would very likely close part of the gap to the
ID-only baseline on warm-start quality.

\paragraph{A production-scale catalogue.} The honest caveat in
Section~\ref{sec:results} is partly an artifact of the synthetic fixture's
small catalogue. Running the same pipeline against the real Kaggle dataset
(\texttt{--kaggle}, $\sim$8k products and $\sim$1M reviews) would both sharpen
the ingredient-affinity signal and lower the chance baseline's competitiveness.

\paragraph{Confidence intervals.} Repeating the cold-start comparison across
several random withheld-population draws and reporting intervals would
establish how much of the zero-versus-positive gap is a reliable structural
property versus a property of one particular random split.

\paragraph{Re-admitting popularity as a tempered prior.} Section
\ref{sec:features} excludes \texttt{item\_popularity\_bucket} from the
retrieval representation entirely because an unweighted inclusion suppressed
cold items. A middle ground -- including it with a separately, more
aggressively regularized embedding dimension, or only for warm items at
inference time -- might recover some of the warm-start ranking quality lost
to the ID-only baseline without reintroducing the cold-item suppression
this report documents.


\bibliographystyle{plain}
\bibliography{references}

\end{document}
